\chapter{Method}
\label{chap:method}

This chapter specifies how the three contributions were built and how they are measured. It covers
the command schema and its grammar (C1), the label-first dataset (C2), the fine-tuning recipe and
quantisation procedure that produce the candidates RQ1 compares, and the evaluation protocol and
definitions of record that the harness of C3 applies. Chapter~\ref{chap:results} reports what that
protocol measured.

\section{Command schema and grammar design}
\label{sec:schema}
The schema is the contract Section~\ref{sec:constrained-decoding} argued for: a fixed vocabulary of
structured commands that a flight controller can act on directly, with every structural property a
controller depends on enforced before a command is dispatched rather than checked afterwards. It
was fixed before the synthetic corpus was generated. No later convenience -- an intent that turned
out to be hard to elicit paraphrases for, a slot that turned out to be hard to fine-tune -- could
then reshape the target the pipeline is built against. It was amended once before generation
completed, to bound the identifier list (below), and every label was checked against the amended
grammar. Ten intents cover the vocabulary: \texttt{formation}, \texttt{move}, \texttt{altitude},
\texttt{takeoff}, \texttt{land}, \texttt{hover}, \texttt{abort}, \texttt{rotate},
\texttt{set\_param}, and \texttt{unknown}. The last is the schema's only provision for an utterance
that does not map onto a command at all.

The wire format is single-line \gls{json} with a key order fixed by the schema declaration and null
or absent slots omitted rather than emitted explicitly:

\begin{lstlisting}[basicstyle=\ttfamily\small, breaklines=true]
{"intent":"formation","shape":"circle","radius":5.0}
{"intent":"move","pos":[10.0,0.0,3.0],"speed":1.5}
{"intent":"move","dir":"north","dist":10.0}
{"intent":"altitude","z":4.0,"ids":[1,2]}
{"intent":"rotate","yaw":90.0}
{"intent":"set_param","speed":1.0}
{"intent":"hover"}
{"intent":"abort"}
{"intent":"unknown"}
\end{lstlisting}

Null fields are omitted to shorten the output the decoder must produce. Over the 200 golden labels,
a command in this wire format averages 16.0 tokens under the Qwen2.5 tokeniser. It would average
19.6 tokens if each intent's absent slots were emitted as nulls, and 53.3 if every slot of the
schema were, as a single flat object for all ten intents would require. Under a per-intent grammar
such as the one below, the first alternative is the relevant one; the second applies to a flat
schema. At the 27.93~tok/s that the
selected configuration sustains on the Raspberry~Pi~5 (Chapter~\ref{chap:results}), those savings
are 0.13~s and 1.33~s of decode time per command, recovered without touching the model.

Three layers stand between the decoder and the flight controller, checking in sequence that a
command is well-formed, within the physical envelope, and legal in the current flight state. This
chapter owns only the first. The second and third, a semantic validator and a state machine, consume
the schema and grammar defined here but belong to the \emph{M\'emoire d'Ing\'enieur}'s architecture
chapter, which discusses their design and their contribution to the safety argument.

\paragraph{The decoding grammar.} The first layer is a context-free grammar in the \gls{gbnf}
formalism, which \texttt{llama.cpp} parses once per request and applies at every decoding step by
masking the decoder's token distribution~\cite{llamacpp}. It is written from the schema and held to
it in two ways: a test suite of valid and adversarial commands, and the conformance check of
Section~\ref{sec:dataset}, which parses this grammar file itself rather than a copy of it. No label
the grammar rejects can therefore enter the dataset:

\begin{lstlisting}[basicstyle=\ttfamily\footnotesize, breaklines=true]
root      ::= cmd
cmd       ::= c-form | c-move | c-alt | c-takeoff | c-land | c-hover | c-abort | c-rot | c-param | c-unk

c-form    ::= "{\"intent\":\"formation\",\"shape\":" shape o-radius o-spacing o-ids "}"
c-move    ::= "{\"intent\":\"move\","  ( kv-pos | kv-dir ) o-speed o-ids "}"
c-alt     ::= "{\"intent\":\"altitude\",\"z\":" num o-ids "}"
c-takeoff ::= "{\"intent\":\"takeoff\""  o-z o-ids "}"
c-land    ::= "{\"intent\":\"land\""     o-ids "}"
c-hover   ::= "{\"intent\":\"hover\""    o-ids "}"
c-abort   ::= "{\"intent\":\"abort\"}"
c-rot     ::= "{\"intent\":\"rotate\",\"yaw\":" num o-ids "}"
c-param   ::= "{\"intent\":\"set_param\"" o-speed o-spacing o-alt "}"
c-unk     ::= "{\"intent\":\"unknown\"}"

shape     ::= "\"circle\"" | "\"line\"" | "\"wedge\"" | "\"grid\"" | "\"column\"" | "\"flock\""
dir       ::= "\"north\"" | "\"south\"" | "\"east\"" | "\"west\"" | "\"up\"" | "\"down\"" | "\"forward\"" | "\"back\"" | "\"left\"" | "\"right\""

kv-pos    ::= "\"pos\":[" num "," num "," num "]"
kv-dir    ::= "\"dir\":" dir ",\"dist\":" num

o-radius  ::= ( ",\"radius\":"  num )?
o-spacing ::= ( ",\"spacing\":" num )?
o-speed   ::= ( ",\"speed\":"   num )?
o-alt     ::= ( ",\"alt\":"     num )?
o-z       ::= ( ",\"z\":"       num )?
o-ids     ::= ( ",\"ids\":[" idlist "]" )?
idlist    ::= [0-4] ( "," [0-4] ){0,4}

num       ::= "-"? [0-9] [0-9]? [0-9]? ( "." [0-9] )?
\end{lstlisting}

The \texttt{num} rule bounds every numeric field to three integer digits and one decimal place.
This makes \texttt{1e999}, a twenty-digit integer, or a bare trailing decimal point structurally
unreachable rather than unlikely. No value of that shape is in the grammar's language, so the
decoder cannot sample toward one, whatever the model's own distribution favours. \gls{gbnf} has no
construct for a numeric \emph{range}, and a range spelt out as digit patterns would make the
decoder emit some admissible value in place of the one spoken, with no record that it had done so.
Range enforcement is therefore left to the semantic validator, which clamps an out-of-range value
and logs the change.
The one-decimal limit fixed here is also why the canonical comparator of
Section~\ref{sec:definitions-of-record} rounds every numeric slot to one decimal before comparison:
the metric matches what the decoder can produce, not a precision the schema never promised.

The \texttt{idlist} rule is bounded in the same way. It restricts a drone identifier to $0$--$4$
and caps the list at five entries, matching the fixed swarm size of $N = 5$. The validator's own
identifier rule, $\mathtt{ids} \subseteq \{0 \dots N{-}1\}$, is the semantic half of the same
constraint. An earlier form of this rule repeated one digit with no bound on the repetitions. In a
preliminary test it let the decoder loop on digits under greedy decoding until the output was
truncated, even for a command as simple as ``swarm land''. Bounding the repetition count in the
grammar makes that failure structurally unreachable rather than rare, the argument the \texttt{num}
rule makes for numeric overflow. Because every rule is bounded, the grammar's language is finite,
and its longest string is 90 characters: a formation command carrying every optional slot at its
widest value and five identifiers. No token is shorter than one character, so the 96-token cap
applied to every grammar-constrained decode exceeds it, and no output under the grammar can be
truncated before it closes.

The design rests on the distinction between a rate that a favourable test set and thorough
fine-tuning can inflate, and a property that holds by construction on any input under that cap.
Section~\ref{sec:constrained-decoding} made that argument in the abstract.
Table~\ref{tab:grammar-ablation} in Chapter~\ref{chap:results} makes it concrete: it decodes the
same fine-tuned models with and without this grammar, and so measures how often they leave the
format when nothing prevents them.

\section{Label-first dataset construction}
\label{sec:dataset}
No public corpus covering this vocabulary was found, so the dataset that trains and evaluates every
model in Chapter~\ref{chap:results} was built rather than assembled from an existing resource, which
is the labour Contribution~C2 names. It was built and sealed under a single governing decision:
generation runs \textbf{label first}. The target \gls{json} is written before any sentence exists.
Natural-language phrasings are then produced for it, up to eight per label for the synthetic corpus,
by deterministic expansion of a committed frame bank and lexicon (\texttt{data/surface\_forms.py})
under a fixed seed. The reverse order, writing sentences and
labelling them afterwards, reintroduces the annotation cost this pipeline exists to remove, and it
produces the label noise that comes with a human resolving ambiguous cases at scale. Writing the
label first makes the label the source of the data rather than a judgement made about it. The
finished dataset is frozen under the version tag \texttt{dataset-v1.0} and documented in a dataset
card (\texttt{data/dataset\_card.md}) that records each component's source, licence and
construction. Table~\ref{tab:dataset} lists every component that this document's experiments and
the \emph{M\'emoire d'Ing\'enieur}'s runtime consume, with the role that motivates each one. That
runtime has two paths from speech to command. On the \textbf{reflex path}, a keyword spotter issues
``swarm hold'' or ``swarm abort'' directly. On the \textbf{parse path}, \gls{vad} endpoints the
utterance, declaring it finished once a stretch of silence follows it; speech recognition then
transcribes it, and the language model parses the transcript.

\begin{table}[htbp]
\centering
\footnotesize
\caption[Dataset components]{Dataset components, at their sizes as generated. \texttt{wake\_pos} and
\texttt{wake\_neg} train and evaluate the keyword spotter on which the \emph{M\'emoire
d'Ing\'enieur} builds its reflex path; every other row feeds an experiment of this document.
\texttt{test\_golden} is one recording session, held over two consecutive days (see below). Roles marked * are reported in the
\emph{M\'emoire d'Ing\'enieur}.}
\label{tab:dataset}
\begin{tabularx}{\textwidth}{@{}l >{\raggedright\arraybackslash}p{2.6cm} X X@{}}
\toprule
Split & Size & Source & Role \\
\midrule
\texttt{train\_synth} & 1{,}940 pairs & Label-first generation; 872 transcripts replaced by round-trip recognition output & \acrshort{lora} fine-tuning (Section~\ref{sec:lora}) \\
\texttt{val\_synth} & 240 pairs & Held-out template families, clean text & Checkpoint selection on \acrshort{em} \\
\texttt{test\_synth} & 240 pairs & Held-out template families, clean text & Second, descriptive split of the multi-model benchmark \\
\texttt{test\_golden} & 200 audio files, 200 transcripts & Author, one recording session over two days; transcripts from the 12 \texttt{test\_synth} families & Multi-model benchmark accuracy of record and confirmatory paired tests; speaker-sensitivity experiment; acoustic-robustness experiment* \\
\texttt{test\_ood} & 150 transcripts & Assistant-style queries, unsupported drone requests and truncated fragments (authored), and recogniser output on noise and silence (12 harvested); text only & False-command rate; the \texttt{unknown} intent's test \\
\texttt{commonvoice} & 300 clips, 263 speakers, 15 accent buckets & Mozilla Common Voice 17.0 English test split~\cite{commonvoice,commonvoice17}, stratified by accent & Multi-speaker baseline of the speaker-sensitivity experiment \\
\texttt{wake\_pos} & 3{,}000 clips (1{,}500 per class) & Three Piper voices~\cite{piper}, 20 renditions each, 25 augmentations per rendition & Keyword-spotter training* \\
\texttt{wake\_neg} & 5{,}040 clips, 3.50~h & LibriSpeech~\cite{librispeech}, Speech Commands~\cite{speechcmd}, 720 near-miss clips & Keyword-spotter training and false-accept measurement* \\
\bottomrule
\end{tabularx}
\end{table}

\paragraph{Synthetic generation and splits.} Diversity in the synthetic corpus is enumerated
explicitly rather than left to sampling temperature. The axes cover the ways a command is spoken:
imperative and polite forms, unit variants (metres, m), numeric forms (five, 5, 5.0),
drone subsetting (all drones, drone two, drones one and three), synonym sets per intent, filler
words, self-corrections, and elliptical commands that drop the verb and keep the slots (``grid, 3.7
spacing''). The \texttt{unknown} intent is taught by 81 of the 1{,}940 training pairs: 20
assistant-style queries, 21 requests for capabilities the schema does not have, such as ``follow
me'' and ``deploy the parachute'', and 40 lexical near-misses that contain a command word without
being a command, such as ``hold that thought''. Some of these
pass through the recognition round trip described below, but every one is \texttt{unknown} in its
clean form, so no training pair maps a supported command made unreadable by recognition error to
\texttt{unknown}. Splits are assigned by \textbf{template family},
not by row: every paraphrase descended from the same label lands in the same split, and the 120
families divide 96/12/12 between training, validation and \texttt{test\_synth}. A row-wise split
would place near-duplicate paraphrases of one label on both sides of the train/test boundary and
inflate every accuracy figure that follows. \texttt{check\_leakage.py} enforces this
partitioning as a gate rather than a report. It checks family and surface-form isolation of the
training split against validation, \texttt{test\_synth} and the golden set, surface-form isolation
of those three against one another, noise-partition isolation and wake-corpus isolation; all four
axes are clean. A second, independent gate, the \textbf{conformance check}, parses every label in
every split against \texttt{cmd.gbnf} (Section~\ref{sec:schema}) before it is accepted. A label the
grammar cannot produce would be an unreachable target that caps the accuracy any model could reach
against it.

\paragraph{Audio synthesis and the golden set.} About half of the training rows keep their clean
transcript. The rest are built by round trip: 872 rows replace the transcript with what
\texttt{whisper.cpp}~\cite{whispercpp} recognises from Piper~\gls{tts}~\cite{piper} audio of it,
mixed with noise at 20, 10 or 5~dB \gls{snr}, and 88 carry injected text perturbations. The training
data therefore contains the recognition errors a model meets at inference rather than only clean
text. The validation and \texttt{test\_synth} splits are clean text with no audio.

The evaluation set of record differs in kind from the synthetic splits in its audio, not in its
text. The single-speaker \textbf{golden set} (\texttt{test\_golden} in Table~\ref{tab:dataset}; the
golden split of Chapter~\ref{chap:results}) is the author reading 200 utterances in a quiet room, in
one session held over two consecutive days, never re-synthesised. Its transcripts were drawn by
the same generator from the 12 template families held out for \texttt{test\_synth}, never from the
validation families that select checkpoints. They cover all ten intents, but only 12 of the 34
(intent, slot-set) combinations the generator produces, one family per combination; one of the 12,
\texttt{set\_param} with altitude and spacing, never occurs in training. Every in-domain exact-match
figure is therefore measured on generator-produced phrasings of 12 command patterns; only
\texttt{test\_ood} probes phrasings outside them.

The protocol also specified two further sessions on a 60-utterance subset: a different room with
real background noise, and the same room on a different day, to capture voice variation across time.
Neither was recorded, and the \emph{M\'emoire d'Ing\'enieur} states the consequence for its runtime
validation as a limitation. The \gls{snr} sweep of the acoustic-robustness experiment, reported in
the \emph{M\'emoire d'Ing\'enieur}, does not depend on them: it is mixed digitally from this
session's clean recording, so noise is the only variable it varies. Holding the speaker fixed is an
internal-validity choice, not an external one. Speaker variation is characterised separately, at
the only speaker-sensitive stage, by the speaker-sensitivity experiment
(Section~\ref{sec:speaker-sensitivity}), which positions the author within a multi-speaker distribution drawn from
\texttt{commonvoice}. The parser consumes text and does not see the speaker, but nothing in this
document establishes that end-to-end accuracy transfers to other speakers.

\paragraph{Noise and annotation.} The noise mixed into training and evaluation audio comes from
ESC-50~\cite{esc50} (environmental sound classes, licensed CC BY-NC 3.0) and DREGON~\cite{dregon}
(rotor noise recorded by microphones mounted on a \gls{uav}, free for personal, educational and
academic use). All noise is mixed by one seeded path, \texttt{data/mix\_noise.py}: DREGON and the
ESC-50 helicopter and engine classes, from the augmentation partition, for the training round trip,
and DREGON alone, from the evaluation partition, for the \gls{snr} sweep. Room impulse responses
from the MIT Acoustical Reverberation Survey~\cite{traer2016} are applied only to the wake corpus,
by \texttt{data/wake\_corpus.py}; the parser's audio is not reverberated. The \gls{snr} is computed on active speech -- the
\gls{rms} over the 20~ms frames within 30~dB of the loudest frame -- against the noise \gls{rms},
rather than over the whole file. Over the whole file, leading silence would set the level and make
the nominal ratio depend on how promptly a speaker starts talking.

Annotation quality is examined, not assumed. The author labelled 50 golden-set utterances, five
per intent, and sealed that pass under SHA-256 on the day recording was completed. The same 50 items
were relabelled after an interval, in a different random order, with neither the sealed pass nor the
generated targets visible. No agreement rate is computed from this exercise: at a short interval, a matching
pair of labels is as easily recall as schema clarity, so a percentage would measure memory rather
than ambiguity. What is reported instead is the list of items the two passes label
\emph{differently}. Recall can only push two passes together, never apart, so a disagreement
indicates either a second reading the schema admits or an annotator slip, whatever the interval;
the count bounds the ambiguities found, not those present. Across the 50 items, at an interval of
4.0~days, that count is zero. This null result is consistent both with a schema that admits no
second reading and with an interval too short for recall to decay. It is reported because omitting
an uninformative outcome would be selective reporting.

\paragraph{The wake corpus.} \texttt{wake\_pos} and \texttt{wake\_neg} are built for the keyword
spotter of the \emph{M\'emoire d'Ing\'enieur}'s reflex path. They are described here because
dataset construction belongs to this document and they reuse its synthesis, noise partitions and
leakage gate, not because this document's experiments consume them. Positive clips were rendered by
three Piper voices, twenty renditions each, with the keyword at a random offset in the clip rather
than centred, and augmented with room impulse responses and noise by the corpus builder
\texttt{data/wake\_corpus.py}. The spotter itself is a classifier head fitted over openWakeWord's
frozen feature front end~\cite{oww}. Negatives are read speech from LibriSpeech~\cite{librispeech},
short command words from Speech Commands~\cite{speechcmd}, and forty near-miss phrases authored by
hand, since neither LibriSpeech nor Speech Commands supplies the confusions that matter to this
spotter: short phrases beginning with ``swarm'' or ``hold''.

Positives are split by base rendition (70/10/20), negatives by speaker, and near-misses by phrase,
so each held-out item differs from every training item along the axis a spotter could otherwise
memorise. Held-out clips are mixed from a disjoint augmentation/evaluation partition, recorded per
clip as \texttt{noise\_partition}: room impulse responses and DREGON split 70/30, and ESC-50 by
fold (folds 1--4 against fold 5). A spotter is therefore never tested under the acoustic
conditions it was augmented with.

\paragraph{The multi-speaker baseline.} The speaker-sensitivity experiment draws 300 clips from the
English test split of Mozilla Common Voice release 17.0~\cite{commonvoice,commonvoice17}, through an
ungated mirror of that release. The clips are stratified across fifteen accent buckets spanning 263
speakers rather than sampled uniformly. The accent-labelled clips of that split are 39\% United
States English, so a uniform draw would answer a question about one accent rather than about
speaker variation.

\section{The LoRA recipe}
\label{sec:lora}
Four base models were fine-tuned with \gls{lora}: Qwen2.5-0.5B-Instruct~\cite{qwen25},
SmolLM2-360M-Instruct~\cite{smollm2} and Llama-3.2-1B-Instruct~\cite{llama32}, spanning the
0.36--1.2~B range RQ1 asks about, and H2O-Danube3-500M-Chat~\cite{danube3} as a parameter-matched
control. The recipe was held identical across all four, so that an accuracy difference
Chapter~\ref{chap:results} reports is attributable to the model under an identical training budget
rather than to how it was trained. Table~\ref{tab:lora} specifies every setting.

\paragraph{The parameter-matched control.} The first three models occupy three parameter counts in
three model families, so every pairwise comparison among them varies scale and pretraining recipe at
once. H2O-Danube3-500M-Chat was added as a control matched in size to Qwen2.5-0.5B-Instruct: 514~M
parameters against 494~M, a difference of 4.0\%, counted from the released checkpoints. It comes
from a fourth family and was trained under the recipe of Table~\ref{tab:lora}. Holding scale
approximately fixed across that pair measures a family difference at one size. It does not remove
family from the size comparison, whose endpoints, SmolLM2-360M and Llama-3.2-1B, are two further
families. What the pair can establish is whether a family difference at fixed size is of the same
order as the difference across the whole range, in which case parameter count alone does not order
these models. Three qualifications bound the comparison. It holds on the reference surface only
(Section~\ref{sec:quantisation-procedure}). It was made under a three-epoch budget at which the
control's validation \gls{em} was still rising. Finally, the control's system prompt was folded into its
user turn, because its chat template has no system role.

\begin{table}[htbp]
\centering
\caption[Fine-tuning recipe]{Fine-tuning recipe, identical across all four base models. Precision is fp16 because the
T4 has no native bf16 support (see below).}
\label{tab:lora}
\begin{tabularx}{\textwidth}{@{}l X@{}}
\toprule
Setting & Value \\
\midrule
Method & \acrshort{lora}~\cite{lora}, not QLoRA \\
Rank, alpha, dropout & $r = 16$, $\alpha = 32$, dropout $0.05$ \\
Target modules & \texttt{q\_proj}, \texttt{k\_proj}, \texttt{v\_proj}, \texttt{o\_proj}, \texttt{gate\_proj}, \texttt{up\_proj}, \texttt{down\_proj} \\
Optimiser & AdamW~\cite{adamw}, learning rate $2\times10^{-4}$, cosine schedule, warmup ratio $0.03$ \\
Epochs, batch & 3 epochs; effective batch 32 (per-device 4 on each of two \acrshort{gpu}s under data parallelism, gradient accumulation 4) \\
Sequence length & 256 tokens \\
Loss & Completion-only masking -- prompt tokens contribute no gradient \\
Precision & fp16 (no native bf16 on the T4; see below) \\
Seed & 42, identical for all four runs \\
Checkpoint selection & Best validation \acrshort{em}, not best validation loss \\
\bottomrule
\end{tabularx}
\end{table}

\paragraph{Adaptation method.} \Gls{lora}~\cite{lora} freezes the base weights and trains
a pair of rank-$r$ update matrices per targeted projection. QLoRA~\cite{qlora} adds one further
step, quantising the frozen base weights during training to fit a larger model into a smaller
training budget. Every model here is small enough to fine-tune without quantising its base weights
on the training hardware, so that step would add a training-time quantisation confound without
recovering any memory the run needs. The rank-16 adapters target every linear projection in the
transformer block, attention and feed-forward alike, rather than attention alone. Adapters on
every linear layer are what QLoRA found necessary to match full fine-tuning~\cite{qlora}, and the
memory a narrower adapter would save is not a constraint on the training hardware. Under the fixed
three-epoch budget, only Llama-3.2-1B-Instruct's validation \gls{em} was unchanged between epochs
two and three (0.900). The other three were still rising at the final epoch: Qwen2.5-0.5B-Instruct
by 0.008 (two of 240 items), SmolLM2-360M-Instruct by 0.025 and H2O-Danube3-500M-Chat by 0.038. No
model is shown to have converged, so a difference between models is a difference under an identical
budget, not a difference between converged models.

\paragraph{Completion-only masking.} The loss is computed only over the tokens the model is asked
to produce -- the \gls{json} command. The prompt tokens that precede it, the one-line instruction
and the transcript alike, contribute no gradient. Training on the prompt as well would reward the
model for reproducing input it is never asked to generate at inference time, diluting the gradient
with tokens that carry no information about whether the command that follows is correct.

\paragraph{Checkpoint selection.} The checkpoint kept from each run is the one with the highest
\gls{em} on \texttt{val\_synth} (Table~\ref{tab:dataset}), not the one with the lowest
validation loss. Cross-entropy loss is computed per token and gives partial credit to a \gls{json}
output that is close in token terms but wrong as a command: a single wrong digit in a coordinate can
cost little loss yet fail \gls{em} entirely. Selecting on loss would therefore optimise a proxy
in place of the metric the deployment depends on. Validation shares the generator with every
in-domain test split, but the golden set is drawn only from test families, so checkpoint selection
never sees a family it is later scored on.

\paragraph{Training precision.} Table~\ref{tab:lora} specifies fp16 rather than bf16, and the
training hardware set this choice. All four models were fine-tuned on Kaggle's free-tier
accelerator, two NVIDIA~T4 \glspl{gpu} of the Turing architecture~\cite{nvidiat4} (compute
capability~7.5~\cite{nvidiacc}). Native bf16 arithmetic arrived with the following architecture,
Ampere (compute capability~8.0)~\cite{nvidiaampere}, so the T4 offers the training run no native
bf16 path, and every run used fp16 instead. The substitution is declared once, here, and not repeated in
Chapter~\ref{chap:results}, because it applied identically to all four models: each was trained
under the same precision, on the same accelerator class, with every other setting in
Table~\ref{tab:lora} held fixed. The ranking among the four models that Chapter~\ref{chap:results}
reports is therefore obtained under one precision applied to all four. Uniform substitution need
not have a uniform effect, however. All four checkpoints are released in bf16, and fp16's narrower
exponent range can affect them differently; this is not ruled out.

\paragraph{Training sessions.} The parameter-matched control was trained in a separate session from
the first three, on the same accelerator class and under the same fp16 substitution, but not in the
same session instance. Training on a \gls{gpu} is not guaranteed to be bit-reproducible across
session instances~\cite{pytorchrepro}, so its adapter is not an output of the run that produced the
other three. Retraining all four together would have replaced the three adapters against which every
quantised result in Chapter~\ref{chap:results} was measured; training the control alone in a second
session was the smaller cost. It qualifies the uniformity claimed above, and Chapter~\ref{chap:discussion} carries it
into the threats to validity.

\section{Quantisation procedure}
\label{sec:quantisation-procedure}
All four fine-tuned adapters were merged into their base weights and converted to \gls{gguf} in
fp16. Three were then quantised with \texttt{llama-quantize}~\cite{llamacpp} to two levels -- Q8\_0
and Q4\_K\_M -- producing six deployable artefacts. The fourth, H2O-Danube3-500M-Chat, was
converted to fp16 \gls{gguf} but not quantised, and is reported on the reference surface defined
below only, for a reason established under chat-template parity rather than chosen in advance. Merging and conversion ran on Kaggle, alongside fine-tuning, because
both need the \texttt{torch}~\cite{pytorch} and \texttt{peft}~\cite{peft} stacks, which are not
installed on the workstation. \texttt{llama-quantize} is native C/C++ with no Python
dependency~\cite{llamacpp}, so quantisation ran locally, against the fp16 \gls{gguf} files that
Kaggle produced.

\paragraph{Reference and deployed surfaces.} The cost of deployment is measured, not assumed, by
holding two surfaces apart. The \textbf{reference surface} is the fine-tuned fp16 model -- base
weights with the \acrshort{lora} adapter applied, not merged -- decoded under
\texttt{transformers}~\cite{transformers} on a Kaggle T4, without a grammar: a reference point, not
a deployed configuration. The \textbf{deployed surface} is one of the six quantised \gls{gguf}
artefacts decoded by \texttt{llama.cpp} under the \gls{gbnf} constraint: the artefact and decoder
that ship. Its accuracy is scored on the workstation and its latency and throughput are timed on the
Raspberry~Pi~5 (Section~\ref{sec:definitions-of-record}). Every deployment accuracy figure in
Chapter~\ref{chap:results} comes from the deployed surface; the four-model family comparison and the
zero-shot baseline, which need H2O-Danube3-500M-Chat or the unadapted models, come from the
reference surface. The reference-minus-deployed difference is the accuracy cost of the deployment
pipeline that the harness of Contribution~C3 measures, and Table~\ref{tab:quantisation-delta}
presents it.

That difference does \emph{not} isolate weight precision. Three things change between the surfaces
at once: the numeric format, the inference runtime and the decoding grammar. The decoding settings
differ as well: a 64-token cap and batched decoding on the reference surface, against a
96-token cap and one sequence at a time on the deployed surface. For Llama-3.2-1B-Instruct, whose
validation \gls{em} tied at epochs two and three, the checkpoint differs too: the reference
surface scored the epoch-3 weights, while the deployed artefacts were merged from the epoch-2
adapter that checkpoint selection kept. The difference is therefore the cost of the deployment
pipeline as a whole. Reporting it as the cost of quantisation would attribute to rounding whatever
the other changes contribute.

One term is separable with the data collected here. Decoding the deployed surface a second time with
the grammar disabled holds precision and runtime fixed and isolates the grammar's effect on the
deployed surface, and Chapter~\ref{chap:results} reports it. Attributing the remainder to precision
and runtime assumes that the grammar's effect is the same at both precisions and that the cap,
batching and checkpoint differences contribute nothing. Separating precision
from runtime would require a further surface, which was not measured. The difference is nonetheless
the quantity a deployment decision needs, since the reference surface is not a configuration that
runs on the board. It is a different quantity from the perplexity and generic-benchmark accuracy
changes the quantisation literature reports (Section~\ref{sec:quantisation}).

\paragraph{Chat-template parity.} Fine-tuning used a Hugging Face tokeniser and its chat template;
inference uses \texttt{llama.cpp}, a separate implementation of the same template. If the two
disagree on a special token, a whitespace convention, or a role marker, the mismatch raises no error
at inference time. The model then decodes against a prompt shaped differently from the one it was
trained on, and accuracy degrades with nothing in the pipeline to flag why. Because the failure is silent,
parity is checked in two halves over ten fixed audit prompts. \texttt{test\_template\_parity.py}
confirms that the \texttt{llama.cpp} tokeniser reproduces the Hugging Face token identifiers for the
Hugging Face rendering. \texttt{check\_parity\_gguf.py} renders each prompt through
\texttt{llama.cpp}'s own template and compares both the rendered text and the token identifiers
with the segmentation the adapter was trained under. The second half was added after the first
deployed-surface sweep and confirmed all three quantised models at Q4\_K\_M. The Q8\_0 artefacts were
not run through it, but their tokeniser and chat-template metadata are identical to their Q4\_K\_M
counterparts' in every \gls{gguf} tokeniser field.

Two differences between the stacks are conventions rather than divergences, and the second half
re-aligns both before comparing. First, \texttt{llama.cpp} omits the leading beginning-of-sequence
piece from the rendered prompt, because it inserts that token at tokenisation
time~\cite{llamacppchat}. Second, Llama-3.2's chat template writes the current date into its system
turn~\cite{llama32card}, so it renders the date of the check rather than the date of training. Each re-alignment is applied only if
it makes the two renderings identical; one that merely narrows a mismatch is discarded and the raw
strings are compared, so neither convention can absorb an unrelated divergence. Without the
re-alignments, the gate would reject Llama-3.2-1B-Instruct over the rendered date alone, although its
prompt is otherwise identical to the training prompt. The date re-alignment accepts one real
difference by design: the model was trained with the date of its training session in the prompt, so
every later serving differs from the training prompt by that date.

\paragraph{Exclusion of the control model.} That policy is what kept H2O-Danube3-500M-Chat out of the
quantised set. Its formatted prompt matches byte for byte under both stacks on all ten audit
prompts, so a comparison of rendered text alone would have passed it. The two tokenisers
nonetheless segment that identical text differently: the word \texttt{drone}, which the prompt
contains, is segmented as \texttt{dr}~+~\texttt{one} by the Hugging Face tokeniser and as
\texttt{dro}~+~\texttt{ne} by \texttt{llama.cpp}. The cause was not isolated. The two sequences are
the same length and detokenise to the same string, so nothing short of a token-level comparison
separates them. A quantised artefact for this model would be served under a segmentation it was
never trained on -- the silent degradation the gate exists to detect -- so the model is reported on
the reference surface only. The other three pass on all ten prompts, in rendered text and token
identifiers, which is the evidence the quantised results of Chapter~\ref{chap:results} rest on.

\section{Evaluation protocol and the definitions of record}
\label{sec:definitions-of-record}
Every accuracy number in either document traces to the functions fixed in this section. The
comparator, the canonical form and the accuracy metrics were fixed before any result existed, so
that no metric could be adjusted after its results were known.

\paragraph{Protocol.} Decoding is greedy for every accuracy figure either document reports -- no
sampling temperature, no beam search. On the deployed surface it also runs under the grammar of
Section~\ref{sec:schema}, one sequence at a time and without a prompt cache. Greedy decoding of one
sequence at a time is what makes a model's output for a given input deterministic. This separates
two kinds of measurement. \Gls{em}, the F1 scores, schema validity, safe-failure rate and
false-command rate take reference text as input. For the selected configuration, the parser's output
on text was identical between the workstation and the Raspberry~Pi~5 on all 200 golden items. These
metrics are therefore scored on the workstation, which changes how long they take to obtain but not
their value. \Gls{crr} passes through speech recognition, whose transcripts differ
between the two machines on two to six of the 200 items per condition, clean audio included, so the \emph{M\'emoire
d'Ing\'enieur} measures it on the board. Latency varies with scheduling, cache state and thermal
throttling on the device that ships, so every timing figure either document reports is measured on
the Raspberry~Pi~5 itself, never inferred from the workstation.

The two kinds consequently report different sample sizes for one configuration: because decoding is
deterministic, $n_{\mathrm{accuracy}}$ is the number of test items in the split being scored, while
$n_{\mathrm{latency}}$ is the number of timed repetitions. The protocol as designed called for 200
repetitions per configuration. The multi-model benchmark timed each Q4\_K\_M configuration over 60
repetitions on the board and did not time Q8\_0 there. The latency experiment of the
\emph{M\'emoire d'Ing\'enieur} timed 226 endpointed segments from the 200 golden utterances on the
parse path, and 78 reflex triggers in each of its idle and loaded conditions.
$n_{\mathrm{latency}}$ is therefore reported per experiment as executed. Reporting one $n$ for both kinds, or the design figure in place of the
executed one, would misstate whichever number it was borrowed from.

\paragraph{Timing protocol.} Every timing run on the Raspberry~Pi~5 pins \texttt{llama.cpp} to the
three cores reserved for inference and refuses to start unless the \texttt{performance} frequency
governor is set. The multi-model benchmark also disables swap, keeps one model resident at a time
and drops the page cache before each configuration loads, so that one configuration's file-cache
warmth cannot flatter the next. Each configuration is warmed up before any trial is scored -- ten
minutes in the multi-model benchmark, five in the latency experiment -- and the warm-up trials are
excluded. The multi-model benchmark's run of record was made with an active cooler fitted. It times
with the server's prompt cache enabled, so that its prefill figure is that of a cached prefix. The
accuracy sweeps and the deployed parser run without it, so the latency experiment's prefill is
uncached.

\paragraph{Software environment.} Quantisation, the deployed surface and every timing run used
\texttt{llama.cpp} build b10863 (commit 88ada91c). Speech recognition used \texttt{whisper.cpp}
commit 52a939a with the \texttt{tiny.en} model~\cite{whisper}, prompted with a fixed 25-word command
vocabulary (\texttt{data/asr.py}), both when the round-trip training rows were built and at run
time. Adapters were trained with \texttt{peft} 0.19.1 on the Kaggle image of each training session.
The session that trained the first three models logged \texttt{transformers} 5.0.0 and PyTorch
2.10.0 (CUDA 12.8); the control's session image was not recorded. \gls{gguf} conversion used the
converter at the head of the \texttt{llama.cpp} repository at the time of each session, which was
not pinned. The system prompt is the single line ``You are a drone swarm command parser. Output only
JSON matching schema.'', with the transcript as the user turn. The multi-model benchmark decodes
with a 96-token cap, a 1{,}024-token context and one sequence at a time; the runtime uses a
512-token context. Speech synthesis uses three Piper voices (en\_US-lessac,
en\_US-ryan and en\_GB-alba, all at medium quality), and generation and noise mixing are seeded
with 42.

\paragraph{Canonical form and comparator.} A single function, $\mathrm{canon}(\cdot)$, is applied to
both the model's prediction and the reference label before either is compared to anything, so that
\gls{em} and \gls{crr} are never computed against two differently normalised
objects. Canonicalisation sets the key order to the schema's declaration order and omits absent or null
fields rather than emitting them. It rounds every float to one decimal place, matching the
one-decimal limit the \texttt{num} rule imposes on what the grammar can produce
(Section~\ref{sec:schema}), deduplicates and sorts \texttt{ids} ascending, and strips all
whitespace outside string literals. Two commands are equal if and only if their canonical forms are
byte-identical. Numeric comparison therefore reduces to a single rule: two numeric slots are equal if
and only if $\mathrm{round}(a,1) = \mathrm{round}(b,1)$.

\Gls{em} and \gls{crr} are, by design, the same comparison applied at two points
in the pipeline. \Gls{em} takes reference text as input, so the model is the only variable under
test and the fine-tuned candidates can be separated cleanly; \gls{crr} takes audio
through the whole pipeline. The comparator and the items are the same: the golden set, scored from
its transcripts and from its audio. The gap between the two therefore measures the cost of the audio
front end, endpointing and speech recognition together, under the rule chosen for scoring an
utterance the endpointer splits. The \emph{M\'emoire d'Ing\'enieur} reports that gap as a finding in its own
right rather than folding it into one end-to-end number. Feeding audio to the reference-text
experiment instead would impose a shared \gls{wer} ceiling on every model, compressing whatever
separates them and confounding any remaining difference between the parser and the recogniser in
front of it.

Table~\ref{tab:metrics} defines every metric this document reports and every accuracy metric the
\emph{M\'emoire d'Ing\'enieur} reuses, together with its latency, formation-accuracy and convergence
definitions; the keyword-spotter, recovery and collision metrics that only
the \emph{M\'emoire d'Ing\'enieur} uses are defined there. This document's results
(Chapter~\ref{chap:results}) draw on \gls{em}, the two F1 scores, schema validity, safe-failure
rate, false-command rate, \gls{wer}, throughput, decode latency, peak resident memory and the
throttle proportion. The remaining rows -- \gls{crr}, latency on the reflex and parse
paths, formation accuracy and convergence -- are computed by and reported in the \emph{M\'emoire
d'Ing\'enieur}. It cites the comparator and the canonical form defined here rather than re-deriving
them. The percentile rule applies to both documents.

% xltabular, not table+tabularx: this table is taller than a text page, and as a float it was
% silently dropped from the PDF entirely (LaTeX reported only "Float too large for page").
% xltabular breaks across pages and repeats the header row.
\begin{xltabular}{\textwidth}{@{}>{\raggedright\arraybackslash}p{3.2cm} X l@{}}
\caption[Metric definitions of record]{Metric definitions of record: the function behind each metric and where it is reported
(section of Chapter~\ref{chap:results}, or \textdagger{} for the \emph{M\'emoire d'Ing\'enieur}).
Every accuracy figure in either document is computed as defined here.}
\label{tab:metrics}\\
\toprule
Metric & Definition & Reported in \\
\midrule
\endfirsthead
\toprule
Metric & Definition & Reported in \\
\midrule
\endhead
\Acrlong{em} (\acrshort{em}) & Fraction of items where $\mathrm{canon}(\mathrm{pred})$ equals $\mathrm{canon}(\mathrm{gold})$; input is reference text & \ref{sec:benchmark}--\ref{sec:grammar-ablation} \\
\Acrlong{crr} (\acrshort{crr}) & The same comparison, with input being audio through the full pipeline & \textdagger \\
Intent macro-F1 & Unweighted mean of per-intent F1 over all ten intents present in the reference, so \texttt{unknown} cannot be inflated by frequency & \ref{sec:benchmark} \\
Slot-F1 & Micro-averaged F1 over the multiset of (key, canonical value) pairs, across all items regardless of intent correctness & \ref{sec:benchmark} \\
Schema validity & Fraction of items whose output parses as a schema-conformant command; under the grammar this is 1.0 by construction, and the grammar ablation is the one place it can move & \ref{sec:benchmark}, \ref{sec:grammar-ablation} \\
Safe-failure rate & Among items where \acrshort{em} is 0, the fraction whose dispatched action is \texttt{unknown} or \texttt{hover}; denominator is failed items only & \ref{sec:benchmark}, \textdagger \\
False-command rate & On the out-of-domain set, the fraction of items whose dispatched action -- the output after the validator, where an invalid output falls back to \texttt{hover} -- is anything other than \texttt{unknown}; \texttt{hover} counts as a false command here & \ref{sec:benchmark} \\
Throughput & Decode tokens summed over decode-seconds summed across the aggregation window, not the mean of per-trial ratios; measured on the three cores reserved for inference (\texttt{taskset -c 1-3}, three \texttt{llama.cpp} threads) & \ref{sec:benchmark}, \ref{sec:selection} \\
Peak resident memory & Maximum resident set size of the language-model process and its children, sampled every 20~ms during the timed run; a lower bound on full-stack memory & \ref{sec:benchmark}, \ref{sec:selection} \\
Throttle proportion & Fraction of timed trials during which the board reports a non-zero throttle flag & \ref{sec:benchmark} \\
\Acrlong{wer} (\acrshort{wer}) & Word-level edits summed over utterances, divided by the total number of reference words (corpus \acrshort{wer}), after the fixed normalisation committed in \texttt{eval/norm.py} & \ref{sec:speaker-sensitivity}, \textdagger \\
Speaker position & The author's \acrshort{wer}, with a bootstrap 95\% confidence interval, located within the distribution of per-bucket \acrshort{wer} over the Common Voice accent buckets & \ref{sec:speaker-sensitivity} \\
Latency & Language-model decode and prefill latency, per trial as reported by the \texttt{llama.cpp} server, reported as p50/p95 alongside throughput. Parse path: $T_0$ = end of speech, $T_1$ = return of the publish call on the runtime's command bus; reflex path: $T_0$ = keyword offset, same $T_1$; both paths also report an onset-anchored figure & \ref{sec:benchmark}, \ref{sec:selection}, \textdagger \\
Percentiles & Nearest-rank on the sorted sample, no interpolation; p50 and p95 denote the 50th and 95th percentiles; sample size $n$ stated in every cell & all \\
Formation accuracy & Fraction of drones within $\tau = 0.5$~m of the assigned slot, averaged over the final 5~s of a 60~s trial & \textdagger \\
Convergence & Rate of trials whose formation accuracy reaches 0.85 at all, plus median and interquartile range of convergence time over converged trials only & \textdagger \\
\bottomrule
\end{xltabular}

\paragraph{Requirements of record.} The metrics above are scored against budgets fixed alongside the
schema before any result existed. Table~\ref{tab:requirements} lists every requirement this document
uses, by the name Chapter~\ref{chap:results} refers to it with, together with its budget and the
experiment that measures it. The throughput floor follows from the decode allowance of
Table~\ref{tab:latency-budget}: a command budgeted at 22 tokens, a design allowance above the
16.0-token mean measured in Section~\ref{sec:schema}, must decode within 1{,}100~ms. Two properties govern how Chapter~\ref{chap:results} may use the list.
Schema validity is satisfied by construction, so it is reported but never counted as a result. The
memory ceiling is defined over the full deployed stack, of which the multi-model benchmark's peak
resident memory measures only the language-model process. That figure can rule a configuration out
but cannot confirm one, and no full-stack measurement was made, so the ceiling is met for the
language-model process only.

\begin{table}[htbp]
\centering
\footnotesize
\caption[Stage latency allowances]{Stage allowances of the parse-path latency budget used as deployment constraints in this
document (p95 targets; the end-to-end row is measured from end of speech). The full budget, including the reflex path,
is derived in the \emph{M\'emoire d'Ing\'enieur}.}
\label{tab:latency-budget}
\begin{tabular}{@{}lr@{}}
\toprule
Stage & Target p95 \\
\midrule
Endpointing (wait to confirm end of speech) & 500~ms \\
Speech recognition (\texttt{whisper.cpp tiny.en}) & 1{,}200~ms \\
Language-model prefill (cached prefix) & 250~ms \\
Language-model decode & 1{,}100~ms \\
Validation, state-machine check and dispatch & 50~ms \\
\midrule
Language-model stages combined & 1{,}350~ms \\
End to end, from end of speech & 2{,}500~ms \\
\bottomrule
\end{tabular}
\end{table}

\begin{table}[htbp]
\centering
\footnotesize
\caption[Requirements of record]{Requirements of record used in this document, with budget and measuring experiment.
Budgets were fixed before any result existed; a missed requirement is reported as missed, not
adjusted. Experiments marked * are reported in the \emph{M\'emoire d'Ing\'enieur}.}
\label{tab:requirements}
\begin{tabularx}{\textwidth}{@{}>{\raggedright\arraybackslash}p{2.3cm} X l >{\raggedright\arraybackslash}p{2.6cm}@{}}
\toprule
Criterion & Requirement & Budget & Measured by \\
\midrule
End-to-end latency & Parse-path latency, from end of speech & p95 $\leq$ 2{,}500~ms & Latency experiment* \\
Exact-match threshold & \acrshort{em} on the deployed quantised artefact, best model & $\geq$ 0.85 & Multi-model benchmark \\
Throughput floor & Decode throughput on the three inference cores, from a 22-token decode allowance within the 1{,}100~ms decode stage & $\geq$ 20~tok/s & Multi-model benchmark \\
Intent F1 & Intent macro-F1, best model & $\geq$ 0.90 & Multi-model benchmark \\
Schema validity & Schema validity under the grammar & 100\%, by construction & Multi-model benchmark \\
Clean-audio recognition & \acrshort{crr} on the golden set, clean audio & $\geq$ 0.80 & Acoustic-robustness experiment* \\
Recognition in noise & \acrshort{crr} at 10~dB \acrshort{snr} & $\geq$ 0.65 & Acoustic-robustness experiment* \\
Safe failure & Safe-failure rate: errors resolving to \texttt{unknown} or \texttt{hover} & $\geq$ 0.70 & Multi-model benchmark, acoustic-robustness experiment* \\
Memory ceiling & Peak resident memory, selected configuration, full stack; only the language-model process was measured & $\leq$ 2.5~GiB & Multi-model benchmark \\
Memory per configuration & Peak resident memory, every configuration & reported; no pass/fail & Multi-model benchmark \\
Throttling & Proportion of timed trials with a non-zero throttle flag & $\leq$ 5\%, reported for every run & Multi-model benchmark \\
Speaker sensitivity & Speech-recognition \acrshort{wer} on the author's speech, positioned against a multi-speaker distribution & bootstrap 95\% confidence interval; no pass/fail & Speaker-sensitivity experiment \\
False commands & False-command rate on the out-of-domain set & $\leq$ 0.05 & Multi-model benchmark \\
\bottomrule
\end{tabularx}
\end{table}

\paragraph{Selection rule and statistical plan.} The deployed configuration is the one with the
highest \gls{em} among configurations satisfying both the 2{,}500~ms end-to-end latency budget and
the 2.5~GiB memory ceiling, ties broken toward lower p95 latency. The rule was fixed in writing
before any result existed. A configuration satisfies a constraint when the quantity the constraint
bounds has been measured for it and is within the bound. A measurement of part of that quantity
cannot show that a configuration satisfies the constraint, but it can show that it does not: the
benchmark times the language-model stages only, and a configuration whose stages alone exceed the
end-to-end budget misses it. If no configuration satisfies both constraints, the rule fails openly:
each unmet constraint is re-baselined against the measured figure, and the re-baselining is
reported rather than the constraint treated as met.
Paired deployed-surface predictions, under the grammar, are compared by McNemar's
test~\cite{mcnemar1947} in two confirmatory families of three, both on the golden split: model
against model at Q4\_K\_M, and Q8\_0 against Q4\_K\_M within each model. Bonferroni correction within
each family gives $\alpha = 0.0167$ against a family-wise $\alpha = 0.05$; the exact binomial test is
used below 25 discordant pairs and the continuity-corrected chi-square from 25 upwards. Two further
comparisons were added after the results existed and are reported as exploratory: each fine-tuned
reference model against its quantised artefacts, corrected over those six pairs, and the
parameter-matched family control. The tests
treat the 200 golden items as independent, whereas they are paraphrases of 12 command patterns, 10
to 22 items each. Items within a pattern share whatever makes that pattern easy or hard, so the
effective sample is smaller than 200, and the intervals and $p$-values are optimistic to that
extent.
The four base models are also scored zero-shot, without fine-tuning and without the grammar, as the
baseline the fine-tune is measured from.

\paragraph{Stage latency allowances.} The stage allowances against which
Chapter~\ref{chap:results} measures the multi-model benchmark are part of an end-to-end latency
budget derived in the \emph{M\'emoire d'Ing\'enieur}. This document uses it only as a constraint,
and Table~\ref{tab:latency-budget} reproduces the rows it needs. The two language-model rows sum to
the 1{,}350~ms stage allowance used in Figure~\ref{fig:pareto}. The allowances are targets, not a
decomposition of the end-to-end budget: they are set per stage, and the measured end-to-end figure,
not their sum, is what the 2{,}500~ms budget is judged against. Chapter~\ref{chap:results} applies
these definitions, requirements and allowances to the measured configurations, in the order the
selection rule needs them.

